\question[3]
Gegeben ist die Funktion für die binäre Suche und eine sortierte Liste:
\begin{lstlisting}
def binaereSuche(a,x,L=0,R=None):
    if R is None: R = len(a)-1
    if L > R: return -1
    mid = L + (R - L) // 2
    if a[mid] == x: return mid
    if x < a[mid]:
        return binaereSuche(a,x,L,mid-1)
    else:
        return binaereSuche(a,x,mid+1,R)

a = [3, 8, 15, 21, 30, 42, 57, 64, 71]
\end{lstlisting}
Gib für die Aufrufe \texttt{binaereSuche(a,64)} und \texttt{binaereSuche(a,10)} jeweils die Werte von
\texttt{L} und \texttt{R} bei allen Aufrufen und den Rückgabewert an.
\begin{solutionbox}{4cm}
\begin{lstlisting}[language={}]
binaereSuche(a,64):  L=0 R=8,  L=5 R=8,  L=7 R=8           Rückgabe 7
binaereSuche(a,10):  L=0 R=8,  L=0 R=3,  L=2 R=3,  L=2 R=1  Rückgabe -1
\end{lstlisting}
\end{solutionbox}
